% Chapter 3

\chapter{Être développeur à Cadremploi} % Main chapter title

\label{apports} % For referencing the chapter elsewhere, use \ref{apports_techniques}

\lhead{Troisième partie. \emph{Être développeur à Cadremploi}} % This is for the header on each page - perhaps a shortened title

%-------------------------------------------------------------------------------
\begin{figure}[h]
  \begin{center}
    \hspace*{-1in}
    \includegraphics[width=0.5\textwidth]{Pictures/CE_logo.png}
  \end{center}
\end{figure}
%-------------------------------------------------------------------------------
\paragraph{}
Le site Cadremploi.fr est découpé en différentes parties qui sont en réalité des applications à part.
La partie publique est celle dans laquelle un utilisateur peut rechercher une offre en fonction de critères et y postuler.
Une seconde partie, destinée aux recruteurs, leur permet de saisir une offre d'emploi qui sera ensuite consultable depuis la partie publique.
J'ai travaillé durant ce stage sur la partie concernant les recruteurs, ou "Espace Recruteur".
Il était en effet à mon arrivée difficilement maintenable, utilisable et son interface était obsolète.
Une refonte totale en était nécessaire.
%-------------------------------------------------------------------------------
\paragraph{}
Mon stage consistait ainsi à participer avec l'équipe Cadremploi à la création d'un nouvel Espace Recruteur, plus moderne.
Je vais détailler dans cette partie le fonctionnement du nouvel espace recruteur ainsi que la façon dont l'équipe s'est organisé pour le développer.


\section{Cahier des charges et fonctionnement général du cycle de développement}
%----------------------------------------------------------------------------------------
\subsection{Contenu du stage}
\label{sub:Contenu du stage}

Le sujet principal de mon stage concernait la refonte de l'espace recruteur du site cadremploi.fr.
CADREMPLOI.fr a conçu un espace réservé aux professionnels afin de les aider dans leurs campagnes de recrutement.
Grâce à ce service de e-commerce, il leur est possible d'accéder à tous les produits et services qu'offre Cadremploi et de payer directement en ligne avec leur carte de crédit ou par chèque après réception de la facture.
%%%%%%%%%%%TODO%%%%%%%%%%%%%%%%%%%%%%
Cette application web avait déjà fait l'objet d'un développement il y a BLABLA ans et une refonte totale en était nécessaire, de manière à proposer des services nouveaux ainsi qu'un design plus actuel.
Concrètement, les principaux objectifs de ce stage étaient les suivants:
\begin{itemize}
  \item{} Développer de nouvelles fonctionnalités dans un environnement technique dynamique (Play! Scala, ElasticSearch, AngularJS)
  \item{} Maintenir le haut niveau de qualité et de tests présents
  \item{} Participer aux ateliers d'architecture et de cadrages techniques
  \item{} Echanger autour de bonnes pratiques avec les équipes de développement
\end{itemize}
%%%%%%%%%%%TODO%%%%%%%%%%%%%%%%%%%%%%
Les fonctionnalités attendues comprenaient notamment la mise en place du pattern Event Sourcing, que je détaillerai plus tard, de manière à pouvoir suivre de manière rigoureuse l'évolution de chaque utilisateur dans l'espace recruteur, ainsi que la simplification du système de classification des offres en créant un produit \"de base\" auquel se grefferont des options.

%----------------------------------------------------------------------------------------
\subsection{Méthodologie de management}
\label{sub:Méthodologie de management}
\paragraph{}
L'équipe Cadremploi, sur le projet Espace Recruteur notamment, suit une méthodologie Agile.
%%%%%%%%%%%TODO%%%%%%%%%%%%%%%%%%%%%%
% Bref résumé de ce qu'est la méthode Agile

% Rappel, MEP chaque semaine + MEPs exceptionnelles, sprint de 2 semaines
Le système de tickets post-it, des stand-up-meeting était très bien suivi bien que leur pratique ne soit pas un exercice figé.
Le contenu des stand-up-meeting est par exemple passé au cours de mon stage de "Qu'est-ce que j'ai fait?, Qu'est-ce que je vais faire?" à des questions plus pratiques: "Qu'est-ce que j'ai fait de réellement impactant pour le reste de l'équipe? Quels dysfonctionnements ai-je remarqué?".
Finalement, plutôt que l'évolution de chacun dans sa tâche, c'est l'évolution de l'application elle-même qui était alors discutée durant les stands-up, réduisant ainsi le temps de l'exercice, ce qui était nécessaire dans cette équipe de 10 personnes.
\paragraph{}
Les post-its suivaient un cycle de vie dont les étapes rentraient dans les catégories suivantes:
\begin{itemize}
  \item{Nouveau}, il s'agit du statut des tickets qui viennent d'arriver sur le tableau et qui seront donc à traiter prochainement
  \item{En cours}, il s'agit d'un ticket qu'un membre de l'équipe est en train de traiter.
  \item{Recette équipe}, lorsqu'une tâche est terminée, elle doit être validée par un membre de l'équipe; ce système de team review que je traiterai plus tard permet d'éviter de nombreuses erreurs au cours du développement de l'application.
  \item{Recette PO} lorsque le ticket passe l'étape de recette équipe, il arrive en recette PO.
  Le Project Owner s'occupe donc de vérifier que les modifications apportées correspondent bien aux attentes et qu'elles ne perturbent pas les anciennes fonctionnalités.
  Si un problème est rencontré, le PO rédige ses retours et la personne responsable du ticket se charge de corriger les problèmes.
  Le ticket reste dans le même état mais un jeu d'étiquettes "Retours Recette" et "Corrigé" est déposé sur le ticket pour qu'il soit possible de suivre l'évolution des corrections.
  \item{Validé}, le PO valide le ticket et les modifications seront mises en production sous peu
  \item{En attente}, le ticket ne peut être traité pour le moment; une synchronisation avec une autre équipe est nécessaire ou bien des questions restent sans réponse.
\end{itemize}
Ainsi les membres de l'équipe choisissent arbitrairement un ticket qu'il traitent en collant une étiquette les représentant dessus et déplacent le ticket dans la catégorie à laquelle il correspond.
\paragraph{}
Le tableau sur lequel sont affiché les tickets sont plutôt grands et plusieurs méthodes sont utilisés pour permettre à l'équipe de s'y retrouver rapidement.
En effet, en plus des colonnes de catégories dans lesquelles sont rangés les tickets, un code couleur permettait de différencier les tâches techniques des tâches fonctionnelles, ainsi que les tâches traitant de parties de Cadremploi autre que l'espace recruteur (projet de refonte de la Home, ...).
De plus, de petites étiquettes, en plus de celles représentant les membres de l'équipe, sont placées sur les tickets et permettent d'indiquer facilement si un des posts-its représente une tâche qui doit être mise en production dans la semaine (Cible MEP), ou si des retours suite à une recette (équipe ou PO) sont à traiter.
\paragraph{}
Enfin, les post-its représentant nos tâches rédigés sur un tableau se trouvent aussi référencés sur l'outil Redmine, une application web de gestion de projets, qui stocke toutes les informations concernant une tâche donnée.
Ce stockage double de l'information est un peu lourd; en effet, il est nécessaire de déplacer le ticket sur le tableau et d'effectuer le même changement en ligne sur Redmine.
Le point positif est qu'il permet à toute l'équipe, même en télétravail, d'accéder facilement à l'information concernant une tâche, notamment aux différentes discussions ayant eu lieu concernant cette tâche.
Il est cependant nécessaire de conserver l'affichage au tableau puisqu'il est bien plus visuel et permet à l'équipe d'échanger bien plus facilement que si le support était uniquement digital.

\section{Cahier des charges et fonctionnement général du cycle de développement}
%-------------------------------------------------------------------------------
\subsection{Contenu du stage}
\label{sub:Contenu du stage}
Le sujet principal de mon stage concernait ainsi la refonte de l'Espace Recruteur du site cadremploi.fr.
L'application actuelle se présente comme un espace réservé aux professionnels afin de les aider dans leur campagnes de recrutement.
Grâce à ce service de e-commerce, il leur est possible d'accéder à tous les produits et services qu'offre Cadremploi et de payer directement en ligne avec leur carte de crédit ou par chèque après réception de la facture.
Cette application web avait déjà fait l'objet d'un développement il y a plusieurs années et une refonte totale en était nécessaire, de manière à proposer des services nouveaux, un design plus actuel et des éléments de maintenabilité plus importants.
Concrètement, les principaux objectifs de ce stage étaient les suivants:
\begin{itemize}
  \item{} Développer de nouvelles fonctionnalités dans un environnement technique dynamique (Play! Scala, ElasticSearch, AngularJS)
  \item{} Maintenir le haut niveau de qualité et de tests présents
  \item{} Participer aux ateliers d'architecture et de cadrages techniques
  \item{} Echanger autour de bonnes pratiques avec les équipes de développement
\end{itemize}
\begin{figure}[h]
  \begin{center}
    \hspace*{-1in}
    \includegraphics[width=1\textwidth]{Pictures/EspaceRecruteur.png}
  \end{center}
\end{figure}

%-------------------------------------------------------------------------------
\subsection{Méthode de management}
\label{sub:Méthode de management}
%-------------------------------------------------------------------------------
\paragraph{}
L'équipe Cadremploi, sur le projet Espace Recruteur notamment, suit une méthodologie Agile.
Les méthodes de type agile se veulent plus pragmatiques que les méthodes traditionnelles.
Elles impliquent au maximum le client et permettent une grande réactivité à ses demandes.
Les méthodes agiles prônent 4 valeurs fondamentales: l'équipe, qui doit être soudée et dans laquelle la communication est fondamentale, le fonctionnement de l'application, la collaboration avec les clients au delà de la négociation contractuelle ainsi que la flexibilité de la structure du logiciel, acceptant ainsi aisément le changement.
Pour l'équipe Cadremploi, cela se caractérisait par des meetings quotidien devant un tableau rassemblant des tâches demandées par le client écrites sur des post-its, une mise en production des nouveautés développées chaque jeudi, la mise en place de tests unitaires et end-to-end ainsi que des revues de code pour chaque tâche effectuée.
%-------------------------------------------------------------------------------
\paragraph{Le stand-up meeting}
Comme dans de nombreuses méthodes Agiles, la communication au sein de l'équipe passe notamment par un stand-up-meeting quotidien ainsi qu'un affichage mural composé de posts-its représentant les tâches relatives à l'évolution de l'application.
Le système de tickets post-it, des stand-up-meeting était très bien suivi bien que leur pratique ne soit pas un exercice figé.
Le contenu des stand-up-meeting est par exemple passé au cours de mon stage de "Qu'est-ce que j'ai fait?, Qu'est-ce que je vais faire?, Ai-je des problèmes?" à des questions plus pratiques: "Qu'est-ce que j'ai fait de réellement impactant pour le reste de l'équipe? Quels dysfonctionnements ai-je remarqué?".
Finalement, plutôt que l'évolution de chacun dans sa tâche, c'est l'évolution de l'application elle-même qui était alors discutée durant les stands-up, réduisant ainsi le œtemps de l'exercice, ce qui était une nécessité dans cette équipe de 10 personnes.
De plus, suite à un amoncellement de tâches trop important qui produisait un confus global lors de la lecture du tableau, il a été mis en place des règles permettant de limiter le nombre de tickets en cours de traitement, permettant une meilleure organisation de l'équipe.
%-------------------------------------------------------------------------------
\paragraph{Cycle de vie d'une tâche}
Les post-its suivaient un cycle de vie dont les étapes rentraient dans les catégories suivantes:
\begin{itemize}
  \item "Nouveau", il s'agit du statut des tickets qui viennent d'arriver sur le tableau et qui seront donc à traiter prochainement
  \item "En cours", il s'agit d'un ticket qu'un membre de l'équipe est en train de traiter.
  \item "Recette" équipe, lorsqu'une tâche est terminée, elle doit être validée par un membre de l'équipe.
  Ce système de team review permet d'éviter de nombreuses erreurs au cours du développement de l'application: une vérification du code ainsi que des tests écrits sont effectués.
  \item "Recette" PO lorsque le ticket passe l'étape de recette équipe, il arrive en recette PO.
  Le Project Owner s'occupe donc de vérifier que les modifications apportées correspondent bien aux attentes et qu'elles ne perturbent pas les anciennes fonctionnalités.
  Si un problème est rencontré, le PO rédige ses retours et la personne responsable du ticket se charge de corriger les problèmes.
  Le ticket reste dans le même état mais un jeu d'étiquettes "Retours Recette" et "Corrigé" est déposé sur le ticket pour qu'il soit possible de suivre l'évolution des corrections.
  \item "Validé", le PO valide le ticket et les modifications seront mises en production sous peu
  \item "En attente", le ticket ne peut être traité pour le moment; une synchronisation avec une autre équipe est nécessaire ou bien des questions restent sans réponse.
\end{itemize}
Ainsi les membres de l'équipe choisissent arbitrairement un ticket qu'il traitent en collant une étiquette, un autocollant les représentant dessus et déplacent le ticket dans la catégorie à laquelle il correspond.
Les tickets validés sont généralement mis en production chaque semaine, le jeudi, à l'exception de quelques mises en productions (MEP) exceptionnelles.
\paragraph{Organisation du tableau}
Le tableau sur lequel sont affichés les tickets est plutôt grand et plusieurs méthodes sont utilisées pour permettre à l'équipe de s'y retrouver rapidement.
En effet, en plus des colonnes de catégories dans lesquelles sont rangés les tickets, un code couleur permettait de différencier les tâches techniques des tâches fonctionnelles, ainsi que les tâches traitant de parties de Cadremploi autre que l'espace recruteur (projet de refonte de la Home, ...).
De plus, de petites étiquettes, en plus de celles représentant les membres de l'équipe, sont placées sur les tickets et permettent d'indiquer facilement si un des posts-its représente une tâche qui doit être mise en production dans la semaine (Cible MEP), ou si des retours suite à une recette (équipe ou PO) sont à traiter.
\paragraph{Référencement des tâches}
Enfin, les post-its représentant nos tâches rédigées sur un tableau se trouvent aussi référencées sur l'outil Redmine, une application web de gestion de projets, qui stocke toutes les informations concernant une tâche donnée.
Ce stockage double de l'information est un peu lourd; en effet, il est nécessaire de déplacer le ticket sur le tableau et d'effectuer le même changement en ligne sur Redmine.
Le point positif est qu'il permet à toute l'équipe, même en télétravail, d'accéder facilement à l'information concernant une tâche, notamment aux différentes discussions ayant eues lieu concernant cette tâche.
Il est cependant nécessaire de conserver l'affichage au tableau puisqu'il est bien plus visuel et permet à l'équipe d'échanger bien plus facilement que si le support était uniquement digital.

%-------------------------------------------------------------------------------
\section{Une équipe nombreuse}
\label{sec:Une équipe nombreuse}
%-------------------------------------------------------------------------------
Cadremploi, à mon arrivée, comptait alors 10 membres dont 9 développeurs.
Il s'agit d'une équipe d'une taille importante pour un projet Agile.
Ainsi, de l'organisation et une bonne coordination étaient nécessaires.

%-------------------------------------------------------------------------------
\subsection{Revue de code}
\label{sub:Revue de code}
Le développement d'applications demande très souvent à une équipe de personnes de travailler ensemble.
Plus l'équipe grandit plus il est compliqué de garantir une coordination correcte et une bonne organisation.
Les équipes d'ingénieur sont spécialement soumises à ce genre de problèmes puisque du code est quotidiennement partagé entre plusieurs personnes au sein d'une entreprise.
La revue de code aide à partager le savoir et les bonnes pratiques.
%-------------------------------------------------------------------------------
\paragraph{}
Le flux de production de l'équipe Cadremploi demande à chaque membre de l'équipe à passer un ticket terminé dans un état "R7 Equipe".
Le ticket doit alors être revu par un autre développeur de l'équipe.
Ce dernier doit analyser les changements fait au code, s'assurer qu'ils n'introduisent pas de récession dans l'application et vérifier que le code produit suit bien les normes de code imposées par l'équipe.
%-------------------------------------------------------------------------------
\paragraph{}
La revue de code par d'autres membres de l'équipe offre de nombreux avantages autres qu'une simple vérification du code.
En effet, il est clair que n'avoir qu'un unique membre de l'équipe responsable d'un point critique sur l'application est dangereux.
La revue de code distribue le savoir au sein de l'équipe et permet d'informer plusieurs membres d'un changement.
De plus, elle stimule les conversations portant sur la structure du code, l'architecture de l'application et permet ainsi à de nouveaux venus, comme je l'étais, de s'adapter rapidement et donc de devenir productif plus rapidement.
%-------------------------------------------------------------------------------
\paragraph{}
Il ne m'est pas possible de réelement comparer le mode de fonctionnement de la revue de code de l'équipe Cadremploi.
En effet, mes projets passés se faisaient dans des équipes plus réduites, et du temps était rarament pris pour la revue de code.
Cela ne nous empêchait pas de suivre des règles évitant la "possession" du code par un membre de l'équipe, mais rien de plus.
%-------------------------------------------------------------------------------
\paragraph{}
La revue de code effectuée avec rigueur se montre extrêmement utile, il s'agit d'une des activités qui m'a permi d'apprendre le plus au sein de l'équipe Cadremploi.
En effet, en favorisant la communication mais aussi les critiques, positives, au sein de l'équipe, elle permet une amélioration continue des méthodes de l'équipe.

%-------------------------------------------------------------------------------
\subsection{Concertation d'équipe}
\label{sub:Concertation d'équipe}
%-------------------------------------------------------------------------------
\subsubsection{L'exemple de l'utilisation de Git}
\label{subs:Utilisation de Git}
Git est un outil de contrôle de version flexible et puissant.
Il permet notamment de mettre en commun du code de manière à ce que plusieurs personnes puissent contribuer à un même projet.
Les options qu'il propose sont nombreuses et il est parfois compliqué de déterminer la meilleure façon de contribuer à un projet sans nuire aux autres contributeurs, la mise en commun de code étant parfois, si ce n'est souvent, source de conflits entre deux versions écrites par différentes personnes.
Notre équipe est constituée de 9 développeurs, et nous ne pouvions échapper à de tels scénarios.
C'est pourquoi quelques semaines après mon arrivée, nous avons convenu avec l'équipe de bonnes pratiques.
Leurs buts étaient multiples, il s'agissait ici de limiter les problèmes de mise en commun de code souvent trop nombreux et bloquants, de simplifier la relecture par les autres membres de l'équipe lors des revues de code mais aussi de clarifier la lecture de l'historique d'évolution de l'application.
%-------------------------------------------------------------------------------
\subsubsection{Les réunions}
\label{subs:Les réunions}
%-------------------------------------------------------------------------------
\paragraph{}
L'utilisation de Git a ainsi été sujet à une concertation de toute l'équipe, mais cela a été le cas pour de nombreux autres points.
J'avais mentionné précédemment que la communication était un point très important dans une méthodologie Agile, et c'est pourquoi des points ponctuels en dehors des stand-up meetings étaient nécessaires.
Il est souvent arrivé que notre équipe se réunisse dans le but de discuter autour de bonnes pratiques au niveau du code, de l'utilisation du fonctionnel qui est nouveau pour la plupart des membres de l'équipe issus d'un univers Java, mais aussi sur des sujets d'architecture et de modélisation.
Ces réunions étaient souvent organisées surtout dans le but de présenter une ligne directrice, mais tout participant était écouté et l'idée principale se trouvait alors renforcée par les discussions qui émergeaient.

\section{Aperçu des technologies et des techniques utilisées}

\subsection{Langages}
\label{sub:Langages}
\subsubsection{Backend}
\label{subs:Backend}
\begin{figure}[h]
  \begin{center}
    \includegraphics[width=0.3\textwidth]{Pictures/play_logo.png}
    \hspace{1in}
    \includegraphics[width=0.15\textwidth]{Pictures/scala_logo.png}
  \end{center}
\end{figure}
\paragraph{Le framework Play}
Le nouvel espace recruteur est géré en backend via le framework Play! dans sa dernière version (2.4.2) et sous le langage Scala.
Il s'agit d'un couple qui permet d'écrire du code rapidement et de façon maintenable.
Play est l'un des framework les plus connus dans le monde Java et offre de nombreux avantages:
\begin{itemize}
  \item Il supporte le rechargement à chaud.
  Il suffit de lancer Play via sa console en mode développement et il prendra automatiquement en compte à chaud les changements effectués sur le code, mais aussi les templates ou le routage.
  Cela contribue grandement au gain de productivité qu'offre Play.
  \item En plus de s’appuyer sur du code à typage statique, Play propose la sécurité du typage à d’autres endroits et notamment sur les templates ou sur le routage des différents contrôleurs.
  Un certain nombre de problèmes sont alors mis en lumière directement à l’étape de la compilation.
  \item Il permet d'exécuter des tests, notamment sur la couche web à plusieurs niveaux:
  On peut par exemple tester un contrôleur en démarrant un serveur web (donc via HTTP) ou, sans le démarrer, en appelant simplement le contrôleur avec le bon contexte, le tout de manière simple et rapide à l’exécution.
  \item Il est sans état et basé sur des entrées sorties non bloquantes et permet ainsi une capacité à monter en charge très intéressante
  \item Il supporte nativement REST, JSON, Websocket entre autres et se présente donc comme un framework moderne
\end{itemize}
Un des points négatifs de ce choix en backend est l'utilisation de fait obligatoire de SBT qui nous a posé de nombreux problèmes.
En effet, cet outil était souvent lent, compliqué d'utilisation et ralentissait souvent le développement notamment via ses résolutions de dépendances interminables.
%-------------------------------------------------------------------------
\paragraph{Scala}
Scala a été le langage choisi pour le développement de l'espace recruteur pour sa flexibilité et sa richesse.
Il s'agit d'un langage permettant d'implémenter des fonctionnalités complexes, notamment liées à l'asynchronisme, tout en maintenant un code clair et de qualité.
Ce langage a aussi permi d'utiliser la puissance d'Akka ainsi que de nombreuses librairies propres à son écosystème et à celui de Java.
\paragraph{Akka}
Cadremploi utilise Akka pour diffuser des événements de manière asynchrone au sein des applications Espace Recruteur et Backoffice.
Akka est aussi utilisé pour programmer l'exécution de scripts récurrents.

%--------------------------------------------------------------
\subsubsection{Frontend}
\label{subs:Frontend}
\begin{figure}[h]
  \begin{center}
    \includegraphics[width=0.4\textwidth]{Pictures/angular_logo.png}
  \end{center}
\end{figure}
Du côté front, l'équipe a utilisé le framework AngularJS ainsi que la librairie D3.js qui permettent des animations dynamiques.
Ce choix de technologie peut être soutenu pour les points suivants:
\begin{itemize}
  \item Angular permet une maintenabilité de l'application aisée, notamment en demandant une structure de type MVC au développeur, ainsi qu'en utilisant du HTML, qui est déclaratif, pour définir l'interface utilisateur.
  \item Il s'agit de plus d'un framework flexible puisqu'il est possible de composer une application en prototypant des composants qui sont de plus facilement testables.
  \item De très nombreuses librairies sont disponibles ce qui permet des possibilités très larges à l'équipe et notamment une utilisation très interactive des formulaires.
  \item L'intégration des tests y est forte ce qui est important pour l'équipe Cadremploi.
\end{itemize}


\paragraph{}
Cette stack reactive assure à l'équipe Cadremploi une consistence ainsi qu'une clareté architecturale.
Elle permet la création d'une application responsive et fiable, donnant aux utilisateurs et aux clients une confiance résidant dans son architecture et son implémentation.

\subsection{Environnements et outils}
\paragraph{IntelliJ}
\label{par:IntelliJ}
\begin{wrapfigure}{l}{0.15\textwidth}
  \vspace{-2.5em}
  \begin{center}
    \includegraphics[width=0.15\textwidth]{Pictures/intellij_logo.png}
  \end{center}
\end{wrapfigure}
Le développement de l'espace recruteur s'est fait sous le framework Play! et IntelliJ Idea était l'IDE utilisé par l'équipe pour développer, puisqu'il offre une bonne intégration de ce framework.
J'avais déjà souvent utilisé cet IDE pour différents projets que j'ai eu à rendre tout au long de ma scolarité à l'EISTI, ainsi son utilisation ne m'a posé aucun problème.
\paragraph{Git}
\label{par:Git}
L'équipe étant nombreuse et l'application étant sujete à grossir avec le temps, un outil de gestion de versions a été utilisé.
Git est l'outil le plus répandu et le plus efficace connu, c'est ce que l'équipe Cadremploi a utilisé.
Il s'agissait aussi d'un outil que j'avais beaucoup utilisé mais de manière très simpliste et j'ai appris à m'en servir d'une toute autre façon, bien plus complète, pendant ce stage.
\paragraph{Jenkins}
\label{par:Jenkins}
\begin{wrapfigure}{rH}{0.25\textwidth}
  \begin{center}
    \includegraphics[width=0.15\textwidth]{Pictures/jenkins_logo.png}
  \end{center}
\end{wrapfigure}
L'équipe Cadremploi utilise aussi Jenkins comme outil d'intégration continue pour déployer les nouveautés sur les différents environnements de production utilisés (recette, intégration), les mises en production et préproductions n'étant pas gérées directement par l'équipe.
Cet outil avait été brièvement présenté lors du cours d'Outils de Développement mais jamais utilisé.
Jenkins est un outil open source permettant de compiler et de tester un projet de manière continue.
Il aide de fait les développeurs à intégrer facilement des changements à une application.
Bien que nous avions été tenté de mettre en place un Jenkins avec mon équipe de Projet de Fin d'Étude, cela ne s'est jamais fait par manque de temps et puisqu'un besoin concret ne s'est jamais présenté.
Je n'ai pas eu de mal à l'utiliser puisqu'une fois configuré, cet outil est réelement simple d'utilisation.
Plusieurs environnements de développement sont utilisés par Cadremploi comme pour les autres équipes de Figaro Classifieds:
\begin{itemize}
  \item{Environnement de développement}, correspond au poste de développeur qui récupère directement les changements via Git.
  \item{Environnement de recette}, où la Project Owner vérifie que les changements apportés fonctionnent effectivement et les valide.
  \item{Environnement de pré-production} qui est un environnement qui ressemble autant que possible à l'environnement de production où sont effectués les derniers tests avant la mise en production
  \item{Environnement de production} qui est la dernière étape, contenant l'application utlisée réelement par les utilisateurs
  \item{Environnement d'intégration}, il s'agit d'un environnement supplémentaire qui est préféré pour les tests fait sur l'infrastructure du projet.
\end{itemize}
J'ai été habitué durant ma scolarité à l'EISTI à utiliser la majorité des outils présentés précédemment puisque j'étais uniquement étranger à Jenkins.
J'ai malgré tout grandement progressé dans leur utilisation et particulièrement pour Git que j'ai redécouvert; le projet auquel je participais differais grandement en taille de ceux auquels j'avais contribué auparavant.
%Potentiellement, Rundeck, Kibana, Puppet.

\section{Architecture de l'application}
\label{sec:Architecture de l'application}
%-------------------------------------------------------------------------------
L'espace recruteur dispose d'une architecture interne assez singulière et je n'avais jamais eu à faire à ce type d'application auparavant.
En effet, l'équipe Cadremploi a mis en place un modèle d'architecture de type CQRS couplé à une gestion d'état basé sur de l'Event Sourcing.
%-------------------------------------------------------------------------------
\paragraph{CQRS: Command Query Responsibility Seggregation}
\label{par:CQRS: Command Query Responsibility Seggregation}
Dans les systèmes de gestion de donnée traditionnels, les commandes, c'est à dire la mise à jour des données, et les requêtes son exécutées sur un seul regroupement d'entitées regroupées dans une unique base de donnée.
CQRS est un modèle d'architecture plutôt récent dont le principe repose, comme son nom l'indique, sur la séparation entre l'écriture et la lecture de l'information.
Nous avons suivi lors du développement de l'espace recruteur ce pattern puisque la séparation des composants de traitement (les "commands") et de restitution (les "queries") de l'information offrait une architecture très intéressante de laquelle nous avons tiré de nombreux bénéfices tels que la suppression du risque d'effets de bord ou l'allègement des classes de service.
%-------------------------------------------------------------------------------
\paragraph{Event Sourcing}
\label{par:Event Sourcing}
La gestion commune de l'état d'un système consiste à enregistrer l'état courant des objets le composant et de reporter chaque changement effectué sur le système en modifiant directement son état, notamment via une mise à jour en base de données.
L'idée fondamentale de l'Event Sourcing est d'assurer que chaque changement appliqué à l'état d'une application peut être capturée dans un objet de type événement.
Ce n'est plus l'état du système qui est enregistré mais les événements qui ont mené le système à l'état dans lequel il est actuellement.
Cette façon de penser permet notamment d'obtenir un log complet de tous les changements effectués, et donc une traçabilité ainsi qu'une aisance de debug importante.
Il est ainsi possible de pouvoir reconstruire tous les états passés de l'application.
%-------------------------------------------------------------------------------
\paragraph{}
Je vais présenter dans cette partie ce type singulier d'architecture qui permet de mettre en place un système extensible et distribuable.
J'expliquerai ensuite la façon dont elle a été mise en place sur la nouvelle version de l'espace recruteur de cadremploi.fr.

%-------------------------------------------------------------------------------
\input{"Chapters/Chapter3/architecture/organisation"}
\input{"Chapters/Chapter3/architecture/implementation"}
%-------------------------------------------------------------------------------

